% Appendix A: API reference

\chapter{API reference}

\label{AppendixA}

\section{\code{fluxlb.core.gpu}}

\begin{table}[h]
\centering\small
\begin{tabular}{@{}p{0.5\linewidth}p{0.44\linewidth}@{}}
\toprule
Function & Purpose \\
\midrule
\code{get\_device()} & CUDA $>$ MPS $>$ CPU, or \code{FLUXLB\_DEVICE} override \\
\code{resolve\_device(x)} & normalise device-like input \\
\code{accumulation\_dtype(device)} & fp64, or fp32 with warning on MPS \\
\code{get\_gpu\_count()}, \code{is\_gpu\_available()} & CUDA presence \\
\code{get\_gpu\_info()}, \code{poll\_gpu\_status()} & per-GPU name and memory dictionaries \\
\code{device\_summary(device)} & one-line host and device description \\
\code{memory\_allocated}, \code{peak\_memory\_allocated}, \code{reset\_peak\_memory} & allocator counters \\
\code{track\_memory(device)} & context manager yielding \code{MemoryReport} \\
\code{clear\_gpu\_cache()} & GC then release cached memory \\
\code{format\_bytes(n)} & bytes to string \\
\code{synchronize(device)}, \code{Timer(device)} & synchronised timing \\
\bottomrule
\end{tabular}
\caption{Public surface of \code{fluxlb.core.gpu}.}
\end{table}

\section{\code{fluxlb.core.lattices}}

\begin{table}[h]
\centering\small
\begin{tabular}{@{}p{0.34\linewidth}p{0.6\linewidth}@{}}
\toprule
Name & Purpose \\
\midrule
\code{Lattice(c, w, opp, cs2, dtype)} & base \code{nn.Module}; validates shapes and
  integer velocities, registers non-persistent buffers, builds \code{shifts} \\
\code{Lattice.Q}, \code{.D}, \code{.isotropy\_order} & class attributes \\
\code{Lattice.c}, \code{.w}, \code{.opp}, \code{.cs2} & buffers (\cref{tab:lattice-contents}) \\
\code{Lattice.shifts} & list of $Q$ integer tuples for \code{torch.roll} \\
\code{D2Q5(dtype)}, \code{D2Q9(dtype)}, \code{D2Q37(dtype)} & 2D lattices \\
\code{D3Q19(dtype)}, \code{D3Q27(dtype)} & 3D lattices \\
\code{D2Q37.R} & Hermite scaling factor $r$; \code{cs2 = 1 / R**2} \\
\bottomrule
\end{tabular}
\caption{Public surface of \code{fluxlb.core.lattices}.}
\end{table}
